\documentclass{article}

\usepackage[final]{neurips_2024}
\usepackage[utf8]{inputenc}
\usepackage[T1]{fontenc}
\usepackage{hyperref}
\usepackage{url}
\usepackage{booktabs}
\usepackage{amsfonts}
\usepackage{nicefrac}
\usepackage{microtype}
\usepackage{xcolor}
\usepackage{amsmath}
\usepackage{graphicx}

\title{DeepFLAIR*: Evaluating U-Net Architectures for EPI-to-FLAIR* Synthesis}

\author{%
  Alexander Shen \\
  Department of Neurology, Cedars-Sinai Medical Center, Los Angeles, CA
}

\begin{document}

\maketitle

\begin{abstract}
FLAIR* is a synthetic MRI contrast generated by combining two separate
acquisitions, and can be generated from a single scan using deep learning.
Current state-of-the-art image synthesis in medical imaging uses a standard
U-Net. We compare three transformer-based U-Net architectures (Attention U-Net,
Swin-UNETR, and UNETR++ with Voxel-Focused Attention) to a baseline U-Net on
a dataset of 89 subjects using a composite loss of MSE, SSIM, and gradient
consistency. All models except UNETR++ achieve comparable test SSIM
(${\approx}0.918$--$0.924$), while UNETR++ produces visibly blurred outputs
due to overly restricted attention windows. Our findings suggest that in medical
image synthesis tasks, dataset size, memory constraints on batch size, and loss
function design are dominant performance factors, consistent with conclusions
drawn in the segmentation literature. Future work should prioritize patch
sampling strategies and loss weighting over architectural novelty.
\end{abstract}

% -----------------------------------------------------------------------
\section{Introduction}
% -----------------------------------------------------------------------

Multiple sclerosis (MS) is a chronic inflammatory disease which presents
lesions in the brain. These lesions can be seen in Magnetic Resonance Imaging
(MRI), and are particularly visible in T2-FLAIR scans, where they appear as
hyperintense (brighter) against white brain matter~\cite{thompson2018diagnosis,filippi2011magnetic}. The Central Vein
Sign (CVS) is a key biomarker which can be used to differentiate MS lesions
from common mimics~\cite{sati2016central,daboul2023effect}. CVS can be most clearly seen in T2*-weighted
magnitude 3D-EPI images, in which veins appear hypointense. By combining the
T2*-weighted and T2-FLAIR through image multiplication, we can generate a
FLAIR* image which demonstrates hyperintense lesions with dark central
veins~\cite{sati2012flair,sati2017rapid}.

A single high quality MRI scan can take 25 minutes~\cite{ozutemiz2023seven}. Movement
during and between scans can lead to artifacts and misregistration, and long
periods in the scanner may not be accessible to all patients. Workflows which
require lengthy scans, combine multiple acquisitions, or involve offline
post-processing can restrict implementation in clinical settings. Emerging work
has offered deep learning solutions through image synthesis tasks such as
super-resolution~\cite{quah2023improving} and synthetic contrasts~\cite{genc2023deepswi,baburyan2025deepflair}.

Producing an accurate FLAIR* from the T2*-weighted magnitude images is
particularly difficult as it requires learning to suppress the hyperintense
cerebrospinal fluid (CSF) to hypointense while maintaining fine details of
the CVS. Regression-based encoder-decoder architectures, such as the U-Net,
are well suited to this task. These architectures provide the depth to learn
complex features while preserving details through skip connections. In this
paper, we compare the performance of four U-Net architectures on the FLAIR*
synthesis task.

% -----------------------------------------------------------------------
\section{Methods}
% -----------------------------------------------------------------------

We used the CAVS-pilot dataset (N=89) with a train/val/test split of 0.7,
0.1, and 0.2 respectively. FLAIR* volumes were generated from T2* 3D-EPI and
T2-FLAIR acquisitions through linear image registration, interpolation, and
image multiplication~\cite{sati2012flair}. Each volume was individually
intensity-normalized to $[0, 1]$, and padded to $(320, 384, 320)$ voxels.

For training, $(64, 64, 64)$ voxel patches were extracted randomly and
augmented with random flipping in the axial and sagittal directions. Gaussian
blurring was applied at three discrete levels (none, weak, medium) with equal
probability. Augmentation was not applied during validation or testing. Patches
were passed into each architecture with subject multiplication to balance
epoch lengths to 32 steps (Table~\ref{tab:training}). During inference,
sliding window patches of size $(64, 64, 64)$ were taken at stride 32, and
blended with Gaussian overlap.

Loss was calculated through a composite objective function combining
mean-squared-error (MSE) and structural similarity (1-SSIM)~\cite{wang2004ssim}.
A gradient consistency term was also included, calculated as the L1 loss of
finite differences in each spatial dimension, averaged across axes. Gradient
clipping (norm $\leq 1.0$) was applied. All loss weights were set equally to
1.0. Models were optimized with Adam ($\text{lr}=10^{-4}$, $\beta_1=0.9$,
$\beta_2=0.999$).

Models were trained until plateauing or reaching 300 epochs. The checkpoint
producing the lowest validation loss was used for evaluation on MSE and SSIM.
All experiments were conducted on an Ubuntu workstation with a 32-core
3.7~GHz CPU, 256~GB memory, and four NVIDIA Quadro RTX 6000 GPUs.

In total, 4 architectures were tested: the baseline U-Net, Attention U-Net,
Swin-UNETR, and UNETR++ with Voxel-Focused Attention. The U-Net and Swin-UNETR
were additionally evaluated with a sigmoid final activation.

\begin{table}[h]
  \caption{Training configuration per architecture.}
  \label{tab:training}
  \centering
  \begin{tabular}{lcccc}
    \toprule
    Architecture & Batch Size & Subject Mult. & Steps/Epoch & Params \\
    \midrule
    U-Net           & 16 & 8  & 32 & 5.6M  \\
    Attention U-Net & 16 & 8  & 32 & 5.7M  \\
    Swin-UNETR      & 2  & 1  & 32 & 15.7M \\
    UNETR++ (VFA)   & 32 & 16 & 32 & 1.2M  \\
    \bottomrule
  \end{tabular}
\end{table}

\subsection{Baseline U-Net}

The baseline U-Net from the original DeepFLAIR* paper~\cite{baburyan2025deepflair} was
replicated in PyTorch. The encoder has four Double Convolutional layers with
BatchNorm and LeakyReLU, and average pooling steps in between. The bottleneck
is a full Double Convolution layer, and the decoder matches the encoder
dimensions with skip connections allowing features to propagate past the
bottleneck. The output goes to a $1{\times}1{\times}1$ convolutional block
followed by a ReLU activation. We also evaluated a final Sigmoid activation.
BatchNorm was used as larger batch sizes are feasible due to the lower memory
footprint. In the original paper, the network was trained for 75 epochs using
deterministic grid patching and fine-tuned with Optuna at batch size 8.

\subsection{Attention U-Net}

The Attention U-Net~\cite{oktay2018attention} introduces attention gates between the
skip connections of the encoder and decoder. These gates learn which encoder
features should be propagated or suppressed. For example, they can allow CVS
structure to pass through while attenuating hyperintense CSF. The gates are
computed as:

\begin{equation}
  W_x(\mathbf{x}) + W_g(\mathbf{g})
  \xrightarrow{\text{ReLU}}
  \psi\,(1{\times}1{\times}1\;\text{conv})
  \xrightarrow{\text{Sigmoid}}
  \alpha
  \quad\Rightarrow\quad
  \mathbf{x} \cdot \alpha
\end{equation}

where $\mathbf{x}$ is the skip connection, $\mathbf{g}$ is the gating signal
from the decoder, and $\alpha$ is the learned spatial attention mask. The
Attention U-Net is otherwise identical to the baseline U-Net.

\subsection{Swin-UNETR}

Pure Vision Transformers (ViTs) compare each voxel to all others, giving
quadratic complexity in the number of voxels. They also lack the inherent
local spatial inductive bias of CNNs, requiring large datasets to learn
spatial relationships from scratch.

The Sliding-window U-Net Transformer (Swin-UNETR)~\cite{hatamizadeh2022swin,liu2021swin} limits
self-attention to local windows of $(2,2,2)$ tokens, while a shifted-window
strategy allows each token to interact with new tokens across layers, building
a global receptive field. We used a feature size of 24. Due to memory
constraints, batch size was limited to 2. BatchNorm is unstable at small batch
sizes, so instance normalization was used throughout~\cite{ulyanov2017instance}.

\subsection{UNETR++ with Voxel-Focused Attention}

UNETR++~\cite{shaker2024unetrpp} addresses ViT quadratic complexity through Efficient
Paired Attention (EPA) blocks. Each EPA block splits computation into parallel
spatial and channel branches, sharing Q and K projections but with separate V:

\begin{itemize}
  \item \textbf{Spatial branch (VFA):} Standard UNETR++ relies on
    dimensionality reduction which risks information loss. By using Voxel
    Focused Attention~\cite{ntanzi2025vfa}, each voxel computes self-attention
    only within its local $(3,3,3)$ neighborhood, reducing complexity from
    $\mathcal{O}(N^2)$ to $\mathcal{O}(27N)$ while preserving representational
    power.
  \item \textbf{Channel branch:} Each channel is treated as a token, computing
    a $C \times C$ attention matrix to capture global feature relationships
    without spatial cost.
\end{itemize}

The outputs of both branches are then summed and projected. The encoder uses a
3-stage hierarchy with patch size 4, producing $16^3$, $8^3$, and $4^3$ token
grids, with a standard U-Net CNN decoder and skip connections. The architecture
has approximately 1.2M parameters, enabling training at batch size 32 alongside
heavier models.

% -----------------------------------------------------------------------
\section{Results}
% -----------------------------------------------------------------------

Table~\ref{tab:results} reports validation loss (the composite training
objective) and test-set metrics for all architectures. Test metrics are
computed on the held-out test set (18 subjects) on full volumes normalised
to $[0,1]$; SSIM is reported as actual structural similarity (higher is
better), and MSE is mean-squared error (lower is better).

\begin{table}[h]
  \caption{Validation loss and test-set performance. Best result per column in \textbf{bold}.}
  \label{tab:results}
  \centering
  \begin{tabular}{lccc}
    \toprule
    Architecture    & Val Loss $\downarrow$ & Test MSE $\downarrow$ & Test SSIM $\uparrow$ \\
    \midrule
    U-Net (ReLU)    & 0.0737          & \textbf{0.000869} & 0.918          \\
    U-Net (Sigmoid) & \textbf{0.0617} & 0.000914          & \textbf{0.924} \\
    Attention U-Net & 0.0750          & 0.000914          & 0.914          \\
    Swin-UNETR      & 0.0637          & 0.000901          & 0.918          \\
    UNETR++ (VFA)   & 0.1014          & 0.001168          & 0.888          \\
    \bottomrule
  \end{tabular}
\end{table}

Training curves are shown in Figure~\ref{fig:curves}. U-Net (Sigmoid)
achieved the lowest validation loss ($0.0617$ at epoch~192), followed by
Swin-UNETR ($0.0637$ at epoch~230), U-Net ReLU ($0.0737$), and Attention
U-Net ($0.0750$). UNETR++ converged to a higher validation loss
(${\approx}0.10$).

\begin{figure}[h]
  \centering
  \includegraphics[width=\linewidth]{report_files/training_curves.png}
  \caption{Validation loss over training epochs (log scale). Dashed vertical
           lines mark each model's best-checkpoint epoch.}
  \label{fig:curves}
\end{figure}

Qualitative results are shown in Figure~\ref{fig:qual}. U-Net, Attention
U-Net, and Swin-UNETR architectures produce visually plausible FLAIR*
predictions with broadly correct CSF suppression, while maintaining some
lesion hyperintensity and CVS appearance. UNETR++ outputs are noticeably
blurred. Difference maps reveal residual errors concentrated around CSF
boundaries and deep white matter.

\begin{figure}[h]
  \centering
  \includegraphics[width=\linewidth]{report_files/05_004_sagittal_gt_pred_diff.png}
  \caption{Sagittal slice comparison. \textit{Top}: ground-truth FLAIR*.
           \textit{Middle}: predicted FLAIR*. \textit{Bottom}: signed
           difference map (Pred $-$ GT); red indicates over-prediction,
           blue indicates under-prediction ($\pm 0.25$ intensity units).}
  \label{fig:qual}
\end{figure}

% -----------------------------------------------------------------------
\section{Discussion}
% -----------------------------------------------------------------------

\subsection{Architecture Comparison}

Across both validation loss and test metrics, the performance gap between
architectures was narrow. Swin-UNETR achieved the second-best validation
loss ($0.0637$) and matched U-Net (ReLU) on test SSIM ($0.918$), despite
having nearly $3\times$ the parameters and requiring $8\times$ smaller batch
sizes due to memory constraints. Attention U-Net, which adds soft attention
gates to the skip connections, did not improve over the plain U-Net on any
metric, suggesting that the attention mechanism did not provide meaningful
benefit. UNETR++ produced the worst validation loss (${\approx}0.10$) and
visibly blurry outputs, which we attribute to the $(3{\times}3{\times}3)$
local VFA window being too small to reason about local intensity structure,
a limitation that matters more for synthesis than segmentation.

These results echo the conclusions of Isensee et al.~\cite{isensee2021nnu},
who showed that a well-configured U-Net consistently outperforms more complex
architectures on medical image segmentation benchmarks. We find that this
holds for synthesis as well: the dominant bottlenecks in this task are dataset
size, memory constraints that limit effective batch size for transformer-based
models, and complex learning tasks such as CSF suppression and CVS
preservation, none of which are addressed by architectural novelty alone.
The most meaningful performance gain came not from a change in architecture
but from switching the final activation from ReLU to Sigmoid, which
constrained outputs to $[0,1]$ and yielded the best validation loss overall
($0.0617$), as shown in Table~\ref{tab:results}.

\subsection{Loss Functions}

The composite loss function for the baseline training used in DeepFLAIR* was
manually tuned~\cite{baburyan2025deepflair}. We trained models with equal weighting to
evaluate the efficacy of each loss component in image-to-image translation.

During training, validation loss was dominated by the SSIM term, while MSE
and gradient loss remained relatively low. This likely contributes to
difficulty learning CSF hypointensities in the ventricles and brainstem, which
require the model to invert the intensity of entire regions. Increasing MSE
weight empirically reduced MSE without worsening SSIM or gradient loss.

L1 loss was also tested to encourage learning of black background voxels. An
early Swin-UNETR experiment with BatchNorm and sigmoid failed to suppress
intensities with L1, while hard sigmoid succeeded (Figure~\ref{fig:activation}).
L1 also appeared to push CSF intensities darker than the FLAIR* ground truth.

\begin{figure}[h]
  \centering
  \includegraphics[width=0.49\linewidth]{report_files/swin-batchnorm-sigmoid-l1.png}
  \hfill
  \includegraphics[width=0.49\linewidth]{report_files/swin-batchnorm-hardsig-l1.png}
  \caption{Early Swin-UNETR outputs with BatchNorm and L1 loss.
           \textit{Left}: sigmoid activation fails to suppress background.
           \textit{Right}: hard sigmoid successfully suppresses background voxels.}
  \label{fig:activation}
\end{figure}

\subsection{Final Activation Functions}

A ReLU activation is used at the output of the baseline U-Net. This
constrains output to positive values but does not cap outputs below 1.0.
During validation, predicted voxels occasionally exceeded 1.0. Empirically,
this did not appear to affect SSIM loss significantly.

During early training, some architectures produced mostly-black outputs
(near-zero activations). Recovery occurred within several epochs, likely due
to gradient signal from skip connections. Kaiming initialization may have
stabilized early training.

Preliminary Swin-UNETR experiments with instance norm and no final activation
successfully trained for 30 epochs with reasonable outputs, suggesting
activation functions may not be strictly necessary. With BatchNorm, sigmoid
struggled to push background voxels to zero even with L1 loss, while hard
sigmoid succeeded.

\subsection{Learning Rate and Training Dynamics}

The Optuna-tuned Adam parameters from~\cite{baburyan2025deepflair}
($\beta_1=0.468$, $\beta_2=0.94$) appeared too aggressive for learning CSF
hypointensities. Models would briefly produce dark CSF patches, spiking the
loss, then immediately revert. Standard Adam ($\beta_1=0.9$) was used for all
reported experiments. No learning rate scheduler was used; manually annealing
the learning rate after plateau detection would likely improve final
performance.

\subsection{Limitations}

The original DeepFLAIR* implementation uses deterministic grid patching at
stride 32, yielding a fixed number of patches per subject. This was not
feasible for Swin-UNETR due to memory constraints and epoch length. Random
patching was used instead, which may reduce spatial coverage consistency but
could improve generalization.

Gaussian blending was used for sliding window inference. A Hann window would
theoretically reduce boundary artifacts but could not be verified within the
scope of this work. Optuna hyperparameter search was not applied to all
architectures due to compute constraints.

\subsection{Future Work}

Several architectures could warrant further investigation, including DADC
layers (Dual Attention and Depthwise Convolutions) for high-frequency texture
preservation, GAN-based training for perceptual sharpness, and diffusion
models with appropriate conditioning for stochastic synthesis. A learning rate
scheduler (e.g., ReduceLROnPlateau) would likely yield further improvement for
all architectures tested.

Beyond architectural changes, improvements to the training pipeline may yield
stronger gains. Lesion-biased patch sampling, which weights patch centers
toward hyperintense regions, could improve sharpness of MS lesions and CVS
structure in synthesis outputs, which are underrepresented in random patches
over the full brain volume. Anatomically-masked loss functions, computing MSE
or SSIM only within brain tissue or lesion regions, could focus the
optimization on clinically relevant areas rather than penalizing background
voxels equally. Better tuning of the composite loss weights, particularly
increasing the MSE term relative to SSIM, may also help the model learn
regional intensity inversions such as CSF suppression.

% -----------------------------------------------------------------------
\section*{Acknowledgments}
% -----------------------------------------------------------------------

Thank you to Inga for providing me her data and help with this assignment, and to Bryan and Sreekanth for helping me get started.
Thank you to Professor Kao for a fantastic class.

% -----------------------------------------------------------------------
\section*{References}
% -----------------------------------------------------------------------

{\small
\bibliographystyle{unsrt}
\bibliography{references}
}

\end{document}
